%=============================================================================
% Experiment 1 appendix: per-case recovery and calibration figures.
%
% Source tree: outputs/finals/experiment-1/per_case/<family>/<case>/{recovery,ecdf}.pdf
% (frozen snapshot; supersedes outputs/figures/model_class/cf/)
%
% Usage (from a main document):
%   \usepackage{graphicx}
%   \usepackage{placeins}      % for \FloatBarrier
%   \renewcommand{\exonefigdir}{<path to>/outputs/finals/experiment-1/per_case}
%   %=============================================================================
% Experiment 1 appendix: per-case recovery and calibration figures.
%
% Source tree: outputs/finals/experiment-1/per_case/<family>/<case>/{recovery,ecdf}.pdf
% (frozen snapshot; supersedes outputs/figures/model_class/cf/)
%
% Usage (from a main document):
%   \usepackage{graphicx}
%   \usepackage{placeins}      % for \FloatBarrier
%   \renewcommand{\exonefigdir}{<path to>/outputs/finals/experiment-1/per_case}
%   %=============================================================================
% Experiment 1 appendix: per-case recovery and calibration figures.
%
% Source tree: outputs/finals/experiment-1/per_case/<family>/<case>/{recovery,ecdf}.pdf
% (frozen snapshot; supersedes outputs/figures/model_class/cf/)
%
% Usage (from a main document):
%   \usepackage{graphicx}
%   \usepackage{placeins}      % for \FloatBarrier
%   \renewcommand{\exonefigdir}{<path to>/outputs/finals/experiment-1/per_case}
%   %=============================================================================
% Experiment 1 appendix: per-case recovery and calibration figures.
%
% Source tree: outputs/finals/experiment-1/per_case/<family>/<case>/{recovery,ecdf}.pdf
% (frozen snapshot; supersedes outputs/figures/model_class/cf/)
%
% Usage (from a main document):
%   \usepackage{graphicx}
%   \usepackage{placeins}      % for \FloatBarrier
%   \renewcommand{\exonefigdir}{<path to>/outputs/finals/experiment-1/per_case}
%   \input{experiment1_figures}
%
% The default \exonefigdir assumes the including file sits in outputs/docs/.
% Compile experiment1_figures_standalone.tex to preview this section alone.
%=============================================================================

% --- Figure root (override with \renewcommand before \input) -----------------
\providecommand{\exonefigdir}{../finals/experiment-1/per_case}

% --- Per-panel include macros ------------------------------------------------
% #1 = family slug (ddm/rdm/cdm)   #2 = family label (DDM/RDM/CDM)
% #3 = case slug                   #4 = case label
\providecommand{\exonegraphic}[1]{%
  \includegraphics[width=\linewidth,height=0.84\textheight,keepaspectratio]{#1}}

\providecommand{\exonerec}[4]{%
  \begin{figure}[htbp]
    \centering
    \exonegraphic{\exonefigdir/#1/#3/recovery.pdf}
    \caption[#2, #4: recovery]{\textbf{#2, #4 design --- coefficient recovery.}
      Posterior mean (points) with 95\% credible intervals against the
      data-generating value, one panel per active cell of the coefficient
      matrix $\mathbf{B}$. Columns are intrinsic parameters, rows are design
      columns (intercept, main effects, interactions), labeled by regressor and
      category; the panel axes render this as \texttt{predictor|category}.
      Inactive cells are drawn as grey N/A tiles. Point color encodes the row
      kind: intercept (purple), main effect (blue), interaction (red).}
    \label{fig:e1-#1-#3-rec}
  \end{figure}}

\providecommand{\exoneecdf}[4]{%
  \begin{figure}[htbp]
    \centering
    \exonegraphic{\exonefigdir/#1/#3/ecdf.pdf}
    \caption[#2, #4: ECDF calibration]{\textbf{#2, #4 design --- ECDF
      calibration.} Rank-ECDF difference (empirical minus uniform) with
      simultaneous 95\% confidence bands, on the same cell layout as the
      corresponding recovery figure (Figure~\ref{fig:e1-#1-#3-rec}). A
      calibrated posterior stays inside the band; excursions above/below
      indicate bias, and a systematically shallow/steep curve indicates
      over-/under-dispersion.}
    \label{fig:e1-#1-#3-ecdf}
  \end{figure}}

% --- Float placement: these panels are large ---------------------------------
\renewcommand{\topfraction}{0.92}
\renewcommand{\bottomfraction}{0.7}
\renewcommand{\textfraction}{0.06}
\renewcommand{\floatpagefraction}{0.6}

%=============================================================================
\section{Experiment 1: per-case recovery and calibration}
\label{app:experiment-1-figures}
%=============================================================================

This appendix collects the per-condition diagnostic panels for the model-class
CogFormer, which amortizes posterior inference jointly over three
sequential-sampling families --- the drift-diffusion model (DDM), the racing
diffusion model (RDM), and the circular diffusion model (CDM) --- within a
shared generalized-linear design space. For every family we report the six
validation design conditions listed in Table~\ref{tab:e1-conditions}, and for
each condition two diagnostics: \emph{coefficient recovery} (posterior mean and
95\% credible interval against ground truth) and \emph{ECDF calibration}
(rank-based simulation-based calibration, plotted as the difference from the
uniform diagonal with simultaneous 95\% bands). This gives
$3 \times 6 \times 2 = 36$ panels, organized by family in
Sections~\ref{sec:e1-ddm}--\ref{sec:e1-cdm}.

Every panel in this appendix comes from a single frozen checkpoint, so the
figures are directly comparable across families and conditions. Two further
diagnostics --- empirical coverage across nominal interval widths, and a
heatmap of NRMSE, posterior contraction and calibration error --- are produced
for each condition but not reproduced here; the numerical values behind the
latter are tabulated per condition in the accompanying \texttt{metrics.csv}
files.

All panels shown here cover the \emph{local} coefficient matrix only. Global
nuisance parameters (the contamination probability $\pi$) are estimated jointly
with the coefficients but are scored and plotted separately, and their panels
are omitted from this appendix. Within a family, the set of active cells varies
by condition: the ``fixed'' conditions hold the dispersion parameters
$s_v, s_\tau$ at their defaults, so those columns disappear, and conditions
without interaction terms drop the corresponding rows. Cell layout is therefore
comparable across the two diagnostics of a given condition, but not across
conditions. Panel axis labels read \texttt{predictor|category}; ``predictor''
there denotes what the text calls a regressor.

\begin{table}[htbp]
  \centering
  \small
  \begin{tabular}{@{}llp{7.6cm}@{}}
    \toprule
    \# & Condition & Design \\
    \midrule
    1 & Intercept-only    & Intercept row only; all intrinsic parameters free,
                            no regressors. \\
    2 & Regressed         & Intercept plus two continuous main-effect
                            regressors $u_1, u_2$; no interaction. \\
    3 & Fixed-regressed   & As \emph{regressed}, but with the dispersion
                            parameters $s_v, s_\tau$ held fixed. \\
    4 & Interaction       & Intercept, both main effects, and the
                            $u_1{:}u_2$ interaction, with hierarchical
                            shedding of parameters across blocks. \\
    5 & Fixed             & Intercept-only with $s_v, s_\tau$ held fixed. \\
    6 & Mixed-predictor   & Interaction design with $u_1$ categorical (at the
                            training-maximum number of categories) and $u_2$
                            continuous, i.e.\ a categorical~$\times$~continuous
                            interaction. \\
    \bottomrule
  \end{tabular}
  \caption{The six validation design conditions, in the order used throughout
    this appendix. Conditions 1--5 are the benchmark cases; condition~6 is the
    mixed categorical/continuous stress case.}
  \label{tab:e1-conditions}
\end{table}

\FloatBarrier
\clearpage

%-----------------------------------------------------------------------------
\subsection{Drift-diffusion model (DDM)}
\label{sec:e1-ddm}
%-----------------------------------------------------------------------------

Intrinsic parameters: drift rate $v$, threshold separation $a$, relative start
point $z$, non-decision time $\tau$, and the dispersions $s_v, s_\tau$.

\exonerec{ddm}{DDM}{intercept_only}{intercept-only}
\exoneecdf{ddm}{DDM}{intercept_only}{intercept-only}
\exonerec{ddm}{DDM}{regressed}{regressed}
\exoneecdf{ddm}{DDM}{regressed}{regressed}
\exonerec{ddm}{DDM}{fixed_regressed}{fixed-regressed}
\exoneecdf{ddm}{DDM}{fixed_regressed}{fixed-regressed}
\exonerec{ddm}{DDM}{interaction}{interaction}
\exoneecdf{ddm}{DDM}{interaction}{interaction}
\exonerec{ddm}{DDM}{fixed}{fixed}
\exoneecdf{ddm}{DDM}{fixed}{fixed}
\exonerec{ddm}{DDM}{mixed_predictor}{mixed-predictor}
\exoneecdf{ddm}{DDM}{mixed_predictor}{mixed-predictor}

\FloatBarrier
\clearpage

%-----------------------------------------------------------------------------
\subsection{Racing diffusion model (RDM)}
\label{sec:e1-rdm}
%-----------------------------------------------------------------------------

Intrinsic parameters: base drift $v$, drift advantage $v_{\mathrm{diff}}$ for
the favored accumulator, threshold $a$, non-decision time $\tau$, and the
dispersions $s_v, s_\tau$.

\exonerec{rdm}{RDM}{intercept_only}{intercept-only}
\exoneecdf{rdm}{RDM}{intercept_only}{intercept-only}
\exonerec{rdm}{RDM}{regressed}{regressed}
\exoneecdf{rdm}{RDM}{regressed}{regressed}
\exonerec{rdm}{RDM}{fixed_regressed}{fixed-regressed}
\exoneecdf{rdm}{RDM}{fixed_regressed}{fixed-regressed}
\exonerec{rdm}{RDM}{interaction}{interaction}
\exoneecdf{rdm}{RDM}{interaction}{interaction}
\exonerec{rdm}{RDM}{fixed}{fixed}
\exoneecdf{rdm}{RDM}{fixed}{fixed}
\exonerec{rdm}{RDM}{mixed_predictor}{mixed-predictor}
\exoneecdf{rdm}{RDM}{mixed_predictor}{mixed-predictor}

\FloatBarrier
\clearpage

%-----------------------------------------------------------------------------
\subsection{Circular diffusion model (CDM)}
\label{sec:e1-cdm}
%-----------------------------------------------------------------------------

Intrinsic parameters: drift magnitude $v$, drift angle $v_\theta$, boundary
radius $a$, non-decision time $\tau$, and the dispersions $s_v, s_\tau$.

\exonerec{cdm}{CDM}{intercept_only}{intercept-only}
\exoneecdf{cdm}{CDM}{intercept_only}{intercept-only}
\exonerec{cdm}{CDM}{regressed}{regressed}
\exoneecdf{cdm}{CDM}{regressed}{regressed}
\exonerec{cdm}{CDM}{fixed_regressed}{fixed-regressed}
\exoneecdf{cdm}{CDM}{fixed_regressed}{fixed-regressed}
\exonerec{cdm}{CDM}{interaction}{interaction}
\exoneecdf{cdm}{CDM}{interaction}{interaction}
\exonerec{cdm}{CDM}{fixed}{fixed}
\exoneecdf{cdm}{CDM}{fixed}{fixed}
\exonerec{cdm}{CDM}{mixed_predictor}{mixed-predictor}
\exoneecdf{cdm}{CDM}{mixed_predictor}{mixed-predictor}

\FloatBarrier

%
% The default \exonefigdir assumes the including file sits in outputs/docs/.
% Compile experiment1_figures_standalone.tex to preview this section alone.
%=============================================================================

% --- Figure root (override with \renewcommand before \input) -----------------
\providecommand{\exonefigdir}{../finals/experiment-1/per_case}

% --- Per-panel include macros ------------------------------------------------
% #1 = family slug (ddm/rdm/cdm)   #2 = family label (DDM/RDM/CDM)
% #3 = case slug                   #4 = case label
\providecommand{\exonegraphic}[1]{%
  \includegraphics[width=\linewidth,height=0.84\textheight,keepaspectratio]{#1}}

\providecommand{\exonerec}[4]{%
  \begin{figure}[htbp]
    \centering
    \exonegraphic{\exonefigdir/#1/#3/recovery.pdf}
    \caption[#2, #4: recovery]{\textbf{#2, #4 design --- coefficient recovery.}
      Posterior mean (points) with 95\% credible intervals against the
      data-generating value, one panel per active cell of the coefficient
      matrix $\mathbf{B}$. Columns are intrinsic parameters, rows are design
      columns (intercept, main effects, interactions), labeled by regressor and
      category; the panel axes render this as \texttt{predictor|category}.
      Inactive cells are drawn as grey N/A tiles. Point color encodes the row
      kind: intercept (purple), main effect (blue), interaction (red).}
    \label{fig:e1-#1-#3-rec}
  \end{figure}}

\providecommand{\exoneecdf}[4]{%
  \begin{figure}[htbp]
    \centering
    \exonegraphic{\exonefigdir/#1/#3/ecdf.pdf}
    \caption[#2, #4: ECDF calibration]{\textbf{#2, #4 design --- ECDF
      calibration.} Rank-ECDF difference (empirical minus uniform) with
      simultaneous 95\% confidence bands, on the same cell layout as the
      corresponding recovery figure (Figure~\ref{fig:e1-#1-#3-rec}). A
      calibrated posterior stays inside the band; excursions above/below
      indicate bias, and a systematically shallow/steep curve indicates
      over-/under-dispersion.}
    \label{fig:e1-#1-#3-ecdf}
  \end{figure}}

% --- Float placement: these panels are large ---------------------------------
\renewcommand{\topfraction}{0.92}
\renewcommand{\bottomfraction}{0.7}
\renewcommand{\textfraction}{0.06}
\renewcommand{\floatpagefraction}{0.6}

%=============================================================================
\section{Experiment 1: per-case recovery and calibration}
\label{app:experiment-1-figures}
%=============================================================================

This appendix collects the per-condition diagnostic panels for the model-class
CogFormer, which amortizes posterior inference jointly over three
sequential-sampling families --- the drift-diffusion model (DDM), the racing
diffusion model (RDM), and the circular diffusion model (CDM) --- within a
shared generalized-linear design space. For every family we report the six
validation design conditions listed in Table~\ref{tab:e1-conditions}, and for
each condition two diagnostics: \emph{coefficient recovery} (posterior mean and
95\% credible interval against ground truth) and \emph{ECDF calibration}
(rank-based simulation-based calibration, plotted as the difference from the
uniform diagonal with simultaneous 95\% bands). This gives
$3 \times 6 \times 2 = 36$ panels, organized by family in
Sections~\ref{sec:e1-ddm}--\ref{sec:e1-cdm}.

Every panel in this appendix comes from a single frozen checkpoint, so the
figures are directly comparable across families and conditions. Two further
diagnostics --- empirical coverage across nominal interval widths, and a
heatmap of NRMSE, posterior contraction and calibration error --- are produced
for each condition but not reproduced here; the numerical values behind the
latter are tabulated per condition in the accompanying \texttt{metrics.csv}
files.

All panels shown here cover the \emph{local} coefficient matrix only. Global
nuisance parameters (the contamination probability $\pi$) are estimated jointly
with the coefficients but are scored and plotted separately, and their panels
are omitted from this appendix. Within a family, the set of active cells varies
by condition: the ``fixed'' conditions hold the dispersion parameters
$s_v, s_\tau$ at their defaults, so those columns disappear, and conditions
without interaction terms drop the corresponding rows. Cell layout is therefore
comparable across the two diagnostics of a given condition, but not across
conditions. Panel axis labels read \texttt{predictor|category}; ``predictor''
there denotes what the text calls a regressor.

\begin{table}[htbp]
  \centering
  \small
  \begin{tabular}{@{}llp{7.6cm}@{}}
    \toprule
    \# & Condition & Design \\
    \midrule
    1 & Intercept-only    & Intercept row only; all intrinsic parameters free,
                            no regressors. \\
    2 & Regressed         & Intercept plus two continuous main-effect
                            regressors $u_1, u_2$; no interaction. \\
    3 & Fixed-regressed   & As \emph{regressed}, but with the dispersion
                            parameters $s_v, s_\tau$ held fixed. \\
    4 & Interaction       & Intercept, both main effects, and the
                            $u_1{:}u_2$ interaction, with hierarchical
                            shedding of parameters across blocks. \\
    5 & Fixed             & Intercept-only with $s_v, s_\tau$ held fixed. \\
    6 & Mixed-predictor   & Interaction design with $u_1$ categorical (at the
                            training-maximum number of categories) and $u_2$
                            continuous, i.e.\ a categorical~$\times$~continuous
                            interaction. \\
    \bottomrule
  \end{tabular}
  \caption{The six validation design conditions, in the order used throughout
    this appendix. Conditions 1--5 are the benchmark cases; condition~6 is the
    mixed categorical/continuous stress case.}
  \label{tab:e1-conditions}
\end{table}

\FloatBarrier
\clearpage

%-----------------------------------------------------------------------------
\subsection{Drift-diffusion model (DDM)}
\label{sec:e1-ddm}
%-----------------------------------------------------------------------------

Intrinsic parameters: drift rate $v$, threshold separation $a$, relative start
point $z$, non-decision time $\tau$, and the dispersions $s_v, s_\tau$.

\exonerec{ddm}{DDM}{intercept_only}{intercept-only}
\exoneecdf{ddm}{DDM}{intercept_only}{intercept-only}
\exonerec{ddm}{DDM}{regressed}{regressed}
\exoneecdf{ddm}{DDM}{regressed}{regressed}
\exonerec{ddm}{DDM}{fixed_regressed}{fixed-regressed}
\exoneecdf{ddm}{DDM}{fixed_regressed}{fixed-regressed}
\exonerec{ddm}{DDM}{interaction}{interaction}
\exoneecdf{ddm}{DDM}{interaction}{interaction}
\exonerec{ddm}{DDM}{fixed}{fixed}
\exoneecdf{ddm}{DDM}{fixed}{fixed}
\exonerec{ddm}{DDM}{mixed_predictor}{mixed-predictor}
\exoneecdf{ddm}{DDM}{mixed_predictor}{mixed-predictor}

\FloatBarrier
\clearpage

%-----------------------------------------------------------------------------
\subsection{Racing diffusion model (RDM)}
\label{sec:e1-rdm}
%-----------------------------------------------------------------------------

Intrinsic parameters: base drift $v$, drift advantage $v_{\mathrm{diff}}$ for
the favored accumulator, threshold $a$, non-decision time $\tau$, and the
dispersions $s_v, s_\tau$.

\exonerec{rdm}{RDM}{intercept_only}{intercept-only}
\exoneecdf{rdm}{RDM}{intercept_only}{intercept-only}
\exonerec{rdm}{RDM}{regressed}{regressed}
\exoneecdf{rdm}{RDM}{regressed}{regressed}
\exonerec{rdm}{RDM}{fixed_regressed}{fixed-regressed}
\exoneecdf{rdm}{RDM}{fixed_regressed}{fixed-regressed}
\exonerec{rdm}{RDM}{interaction}{interaction}
\exoneecdf{rdm}{RDM}{interaction}{interaction}
\exonerec{rdm}{RDM}{fixed}{fixed}
\exoneecdf{rdm}{RDM}{fixed}{fixed}
\exonerec{rdm}{RDM}{mixed_predictor}{mixed-predictor}
\exoneecdf{rdm}{RDM}{mixed_predictor}{mixed-predictor}

\FloatBarrier
\clearpage

%-----------------------------------------------------------------------------
\subsection{Circular diffusion model (CDM)}
\label{sec:e1-cdm}
%-----------------------------------------------------------------------------

Intrinsic parameters: drift magnitude $v$, drift angle $v_\theta$, boundary
radius $a$, non-decision time $\tau$, and the dispersions $s_v, s_\tau$.

\exonerec{cdm}{CDM}{intercept_only}{intercept-only}
\exoneecdf{cdm}{CDM}{intercept_only}{intercept-only}
\exonerec{cdm}{CDM}{regressed}{regressed}
\exoneecdf{cdm}{CDM}{regressed}{regressed}
\exonerec{cdm}{CDM}{fixed_regressed}{fixed-regressed}
\exoneecdf{cdm}{CDM}{fixed_regressed}{fixed-regressed}
\exonerec{cdm}{CDM}{interaction}{interaction}
\exoneecdf{cdm}{CDM}{interaction}{interaction}
\exonerec{cdm}{CDM}{fixed}{fixed}
\exoneecdf{cdm}{CDM}{fixed}{fixed}
\exonerec{cdm}{CDM}{mixed_predictor}{mixed-predictor}
\exoneecdf{cdm}{CDM}{mixed_predictor}{mixed-predictor}

\FloatBarrier

%
% The default \exonefigdir assumes the including file sits in outputs/docs/.
% Compile experiment1_figures_standalone.tex to preview this section alone.
%=============================================================================

% --- Figure root (override with \renewcommand before \input) -----------------
\providecommand{\exonefigdir}{../finals/experiment-1/per_case}

% --- Per-panel include macros ------------------------------------------------
% #1 = family slug (ddm/rdm/cdm)   #2 = family label (DDM/RDM/CDM)
% #3 = case slug                   #4 = case label
\providecommand{\exonegraphic}[1]{%
  \includegraphics[width=\linewidth,height=0.84\textheight,keepaspectratio]{#1}}

\providecommand{\exonerec}[4]{%
  \begin{figure}[htbp]
    \centering
    \exonegraphic{\exonefigdir/#1/#3/recovery.pdf}
    \caption[#2, #4: recovery]{\textbf{#2, #4 design --- coefficient recovery.}
      Posterior mean (points) with 95\% credible intervals against the
      data-generating value, one panel per active cell of the coefficient
      matrix $\mathbf{B}$. Columns are intrinsic parameters, rows are design
      columns (intercept, main effects, interactions), labeled by regressor and
      category; the panel axes render this as \texttt{predictor|category}.
      Inactive cells are drawn as grey N/A tiles. Point color encodes the row
      kind: intercept (purple), main effect (blue), interaction (red).}
    \label{fig:e1-#1-#3-rec}
  \end{figure}}

\providecommand{\exoneecdf}[4]{%
  \begin{figure}[htbp]
    \centering
    \exonegraphic{\exonefigdir/#1/#3/ecdf.pdf}
    \caption[#2, #4: ECDF calibration]{\textbf{#2, #4 design --- ECDF
      calibration.} Rank-ECDF difference (empirical minus uniform) with
      simultaneous 95\% confidence bands, on the same cell layout as the
      corresponding recovery figure (Figure~\ref{fig:e1-#1-#3-rec}). A
      calibrated posterior stays inside the band; excursions above/below
      indicate bias, and a systematically shallow/steep curve indicates
      over-/under-dispersion.}
    \label{fig:e1-#1-#3-ecdf}
  \end{figure}}

% --- Float placement: these panels are large ---------------------------------
\renewcommand{\topfraction}{0.92}
\renewcommand{\bottomfraction}{0.7}
\renewcommand{\textfraction}{0.06}
\renewcommand{\floatpagefraction}{0.6}

%=============================================================================
\section{Experiment 1: per-case recovery and calibration}
\label{app:experiment-1-figures}
%=============================================================================

This appendix collects the per-condition diagnostic panels for the model-class
CogFormer, which amortizes posterior inference jointly over three
sequential-sampling families --- the drift-diffusion model (DDM), the racing
diffusion model (RDM), and the circular diffusion model (CDM) --- within a
shared generalized-linear design space. For every family we report the six
validation design conditions listed in Table~\ref{tab:e1-conditions}, and for
each condition two diagnostics: \emph{coefficient recovery} (posterior mean and
95\% credible interval against ground truth) and \emph{ECDF calibration}
(rank-based simulation-based calibration, plotted as the difference from the
uniform diagonal with simultaneous 95\% bands). This gives
$3 \times 6 \times 2 = 36$ panels, organized by family in
Sections~\ref{sec:e1-ddm}--\ref{sec:e1-cdm}.

Every panel in this appendix comes from a single frozen checkpoint, so the
figures are directly comparable across families and conditions. Two further
diagnostics --- empirical coverage across nominal interval widths, and a
heatmap of NRMSE, posterior contraction and calibration error --- are produced
for each condition but not reproduced here; the numerical values behind the
latter are tabulated per condition in the accompanying \texttt{metrics.csv}
files.

All panels shown here cover the \emph{local} coefficient matrix only. Global
nuisance parameters (the contamination probability $\pi$) are estimated jointly
with the coefficients but are scored and plotted separately, and their panels
are omitted from this appendix. Within a family, the set of active cells varies
by condition: the ``fixed'' conditions hold the dispersion parameters
$s_v, s_\tau$ at their defaults, so those columns disappear, and conditions
without interaction terms drop the corresponding rows. Cell layout is therefore
comparable across the two diagnostics of a given condition, but not across
conditions. Panel axis labels read \texttt{predictor|category}; ``predictor''
there denotes what the text calls a regressor.

\begin{table}[htbp]
  \centering
  \small
  \begin{tabular}{@{}llp{7.6cm}@{}}
    \toprule
    \# & Condition & Design \\
    \midrule
    1 & Intercept-only    & Intercept row only; all intrinsic parameters free,
                            no regressors. \\
    2 & Regressed         & Intercept plus two continuous main-effect
                            regressors $u_1, u_2$; no interaction. \\
    3 & Fixed-regressed   & As \emph{regressed}, but with the dispersion
                            parameters $s_v, s_\tau$ held fixed. \\
    4 & Interaction       & Intercept, both main effects, and the
                            $u_1{:}u_2$ interaction, with hierarchical
                            shedding of parameters across blocks. \\
    5 & Fixed             & Intercept-only with $s_v, s_\tau$ held fixed. \\
    6 & Mixed-predictor   & Interaction design with $u_1$ categorical (at the
                            training-maximum number of categories) and $u_2$
                            continuous, i.e.\ a categorical~$\times$~continuous
                            interaction. \\
    \bottomrule
  \end{tabular}
  \caption{The six validation design conditions, in the order used throughout
    this appendix. Conditions 1--5 are the benchmark cases; condition~6 is the
    mixed categorical/continuous stress case.}
  \label{tab:e1-conditions}
\end{table}

\FloatBarrier
\clearpage

%-----------------------------------------------------------------------------
\subsection{Drift-diffusion model (DDM)}
\label{sec:e1-ddm}
%-----------------------------------------------------------------------------

Intrinsic parameters: drift rate $v$, threshold separation $a$, relative start
point $z$, non-decision time $\tau$, and the dispersions $s_v, s_\tau$.

\exonerec{ddm}{DDM}{intercept_only}{intercept-only}
\exoneecdf{ddm}{DDM}{intercept_only}{intercept-only}
\exonerec{ddm}{DDM}{regressed}{regressed}
\exoneecdf{ddm}{DDM}{regressed}{regressed}
\exonerec{ddm}{DDM}{fixed_regressed}{fixed-regressed}
\exoneecdf{ddm}{DDM}{fixed_regressed}{fixed-regressed}
\exonerec{ddm}{DDM}{interaction}{interaction}
\exoneecdf{ddm}{DDM}{interaction}{interaction}
\exonerec{ddm}{DDM}{fixed}{fixed}
\exoneecdf{ddm}{DDM}{fixed}{fixed}
\exonerec{ddm}{DDM}{mixed_predictor}{mixed-predictor}
\exoneecdf{ddm}{DDM}{mixed_predictor}{mixed-predictor}

\FloatBarrier
\clearpage

%-----------------------------------------------------------------------------
\subsection{Racing diffusion model (RDM)}
\label{sec:e1-rdm}
%-----------------------------------------------------------------------------

Intrinsic parameters: base drift $v$, drift advantage $v_{\mathrm{diff}}$ for
the favored accumulator, threshold $a$, non-decision time $\tau$, and the
dispersions $s_v, s_\tau$.

\exonerec{rdm}{RDM}{intercept_only}{intercept-only}
\exoneecdf{rdm}{RDM}{intercept_only}{intercept-only}
\exonerec{rdm}{RDM}{regressed}{regressed}
\exoneecdf{rdm}{RDM}{regressed}{regressed}
\exonerec{rdm}{RDM}{fixed_regressed}{fixed-regressed}
\exoneecdf{rdm}{RDM}{fixed_regressed}{fixed-regressed}
\exonerec{rdm}{RDM}{interaction}{interaction}
\exoneecdf{rdm}{RDM}{interaction}{interaction}
\exonerec{rdm}{RDM}{fixed}{fixed}
\exoneecdf{rdm}{RDM}{fixed}{fixed}
\exonerec{rdm}{RDM}{mixed_predictor}{mixed-predictor}
\exoneecdf{rdm}{RDM}{mixed_predictor}{mixed-predictor}

\FloatBarrier
\clearpage

%-----------------------------------------------------------------------------
\subsection{Circular diffusion model (CDM)}
\label{sec:e1-cdm}
%-----------------------------------------------------------------------------

Intrinsic parameters: drift magnitude $v$, drift angle $v_\theta$, boundary
radius $a$, non-decision time $\tau$, and the dispersions $s_v, s_\tau$.

\exonerec{cdm}{CDM}{intercept_only}{intercept-only}
\exoneecdf{cdm}{CDM}{intercept_only}{intercept-only}
\exonerec{cdm}{CDM}{regressed}{regressed}
\exoneecdf{cdm}{CDM}{regressed}{regressed}
\exonerec{cdm}{CDM}{fixed_regressed}{fixed-regressed}
\exoneecdf{cdm}{CDM}{fixed_regressed}{fixed-regressed}
\exonerec{cdm}{CDM}{interaction}{interaction}
\exoneecdf{cdm}{CDM}{interaction}{interaction}
\exonerec{cdm}{CDM}{fixed}{fixed}
\exoneecdf{cdm}{CDM}{fixed}{fixed}
\exonerec{cdm}{CDM}{mixed_predictor}{mixed-predictor}
\exoneecdf{cdm}{CDM}{mixed_predictor}{mixed-predictor}

\FloatBarrier

%
% The default \exonefigdir assumes the including file sits in outputs/docs/.
% Compile experiment1_figures_standalone.tex to preview this section alone.
%=============================================================================

% --- Figure root (override with \renewcommand before \input) -----------------
\providecommand{\exonefigdir}{../finals/experiment-1/per_case}

% --- Per-panel include macros ------------------------------------------------
% #1 = family slug (ddm/rdm/cdm)   #2 = family label (DDM/RDM/CDM)
% #3 = case slug                   #4 = case label
\providecommand{\exonegraphic}[1]{%
  \includegraphics[width=\linewidth,height=0.84\textheight,keepaspectratio]{#1}}

\providecommand{\exonerec}[4]{%
  \begin{figure}[htbp]
    \centering
    \exonegraphic{\exonefigdir/#1/#3/recovery.pdf}
    \caption[#2, #4: recovery]{\textbf{#2, #4 design --- coefficient recovery.}
      Posterior mean (points) with 95\% credible intervals against the
      data-generating value, one panel per active cell of the coefficient
      matrix $\mathbf{B}$. Columns are intrinsic parameters, rows are design
      columns (intercept, main effects, interactions), labeled by regressor and
      category; the panel axes render this as \texttt{predictor|category}.
      Inactive cells are drawn as grey N/A tiles. Point color encodes the row
      kind: intercept (purple), main effect (blue), interaction (red).}
    \label{fig:e1-#1-#3-rec}
  \end{figure}}

\providecommand{\exoneecdf}[4]{%
  \begin{figure}[htbp]
    \centering
    \exonegraphic{\exonefigdir/#1/#3/ecdf.pdf}
    \caption[#2, #4: ECDF calibration]{\textbf{#2, #4 design --- ECDF
      calibration.} Rank-ECDF difference (empirical minus uniform) with
      simultaneous 95\% confidence bands, on the same cell layout as the
      corresponding recovery figure (Figure~\ref{fig:e1-#1-#3-rec}). A
      calibrated posterior stays inside the band; excursions above/below
      indicate bias, and a systematically shallow/steep curve indicates
      over-/under-dispersion.}
    \label{fig:e1-#1-#3-ecdf}
  \end{figure}}

% --- Float placement: these panels are large ---------------------------------
\renewcommand{\topfraction}{0.92}
\renewcommand{\bottomfraction}{0.7}
\renewcommand{\textfraction}{0.06}
\renewcommand{\floatpagefraction}{0.6}

%=============================================================================
\section{Experiment 1: per-case recovery and calibration}
\label{app:experiment-1-figures}
%=============================================================================

This appendix collects the per-condition diagnostic panels for the model-class
CogFormer, which amortizes posterior inference jointly over three
sequential-sampling families --- the drift-diffusion model (DDM), the racing
diffusion model (RDM), and the circular diffusion model (CDM) --- within a
shared generalized-linear design space. For every family we report the six
validation design conditions listed in Table~\ref{tab:e1-conditions}, and for
each condition two diagnostics: \emph{coefficient recovery} (posterior mean and
95\% credible interval against ground truth) and \emph{ECDF calibration}
(rank-based simulation-based calibration, plotted as the difference from the
uniform diagonal with simultaneous 95\% bands). This gives
$3 \times 6 \times 2 = 36$ panels, organized by family in
Sections~\ref{sec:e1-ddm}--\ref{sec:e1-cdm}.

Every panel in this appendix comes from a single frozen checkpoint, so the
figures are directly comparable across families and conditions. Two further
diagnostics --- empirical coverage across nominal interval widths, and a
heatmap of NRMSE, posterior contraction and calibration error --- are produced
for each condition but not reproduced here; the numerical values behind the
latter are tabulated per condition in the accompanying \texttt{metrics.csv}
files.

All panels shown here cover the \emph{local} coefficient matrix only. Global
nuisance parameters (the contamination probability $\pi$) are estimated jointly
with the coefficients but are scored and plotted separately, and their panels
are omitted from this appendix. Within a family, the set of active cells varies
by condition: the ``fixed'' conditions hold the dispersion parameters
$s_v, s_\tau$ at their defaults, so those columns disappear, and conditions
without interaction terms drop the corresponding rows. Cell layout is therefore
comparable across the two diagnostics of a given condition, but not across
conditions. Panel axis labels read \texttt{predictor|category}; ``predictor''
there denotes what the text calls a regressor.

\begin{table}[htbp]
  \centering
  \small
  \begin{tabular}{@{}llp{7.6cm}@{}}
    \toprule
    \# & Condition & Design \\
    \midrule
    1 & Intercept-only    & Intercept row only; all intrinsic parameters free,
                            no regressors. \\
    2 & Regressed         & Intercept plus two continuous main-effect
                            regressors $u_1, u_2$; no interaction. \\
    3 & Fixed-regressed   & As \emph{regressed}, but with the dispersion
                            parameters $s_v, s_\tau$ held fixed. \\
    4 & Interaction       & Intercept, both main effects, and the
                            $u_1{:}u_2$ interaction, with hierarchical
                            shedding of parameters across blocks. \\
    5 & Fixed             & Intercept-only with $s_v, s_\tau$ held fixed. \\
    6 & Mixed-predictor   & Interaction design with $u_1$ categorical (at the
                            training-maximum number of categories) and $u_2$
                            continuous, i.e.\ a categorical~$\times$~continuous
                            interaction. \\
    \bottomrule
  \end{tabular}
  \caption{The six validation design conditions, in the order used throughout
    this appendix. Conditions 1--5 are the benchmark cases; condition~6 is the
    mixed categorical/continuous stress case.}
  \label{tab:e1-conditions}
\end{table}

\FloatBarrier
\clearpage

%-----------------------------------------------------------------------------
\subsection{Drift-diffusion model (DDM)}
\label{sec:e1-ddm}
%-----------------------------------------------------------------------------

Intrinsic parameters: drift rate $v$, threshold separation $a$, relative start
point $z$, non-decision time $\tau$, and the dispersions $s_v, s_\tau$.

\exonerec{ddm}{DDM}{intercept_only}{intercept-only}
\exoneecdf{ddm}{DDM}{intercept_only}{intercept-only}
\exonerec{ddm}{DDM}{regressed}{regressed}
\exoneecdf{ddm}{DDM}{regressed}{regressed}
\exonerec{ddm}{DDM}{fixed_regressed}{fixed-regressed}
\exoneecdf{ddm}{DDM}{fixed_regressed}{fixed-regressed}
\exonerec{ddm}{DDM}{interaction}{interaction}
\exoneecdf{ddm}{DDM}{interaction}{interaction}
\exonerec{ddm}{DDM}{fixed}{fixed}
\exoneecdf{ddm}{DDM}{fixed}{fixed}
\exonerec{ddm}{DDM}{mixed_predictor}{mixed-predictor}
\exoneecdf{ddm}{DDM}{mixed_predictor}{mixed-predictor}

\FloatBarrier
\clearpage

%-----------------------------------------------------------------------------
\subsection{Racing diffusion model (RDM)}
\label{sec:e1-rdm}
%-----------------------------------------------------------------------------

Intrinsic parameters: base drift $v$, drift advantage $v_{\mathrm{diff}}$ for
the favored accumulator, threshold $a$, non-decision time $\tau$, and the
dispersions $s_v, s_\tau$.

\exonerec{rdm}{RDM}{intercept_only}{intercept-only}
\exoneecdf{rdm}{RDM}{intercept_only}{intercept-only}
\exonerec{rdm}{RDM}{regressed}{regressed}
\exoneecdf{rdm}{RDM}{regressed}{regressed}
\exonerec{rdm}{RDM}{fixed_regressed}{fixed-regressed}
\exoneecdf{rdm}{RDM}{fixed_regressed}{fixed-regressed}
\exonerec{rdm}{RDM}{interaction}{interaction}
\exoneecdf{rdm}{RDM}{interaction}{interaction}
\exonerec{rdm}{RDM}{fixed}{fixed}
\exoneecdf{rdm}{RDM}{fixed}{fixed}
\exonerec{rdm}{RDM}{mixed_predictor}{mixed-predictor}
\exoneecdf{rdm}{RDM}{mixed_predictor}{mixed-predictor}

\FloatBarrier
\clearpage

%-----------------------------------------------------------------------------
\subsection{Circular diffusion model (CDM)}
\label{sec:e1-cdm}
%-----------------------------------------------------------------------------

Intrinsic parameters: drift magnitude $v$, drift angle $v_\theta$, boundary
radius $a$, non-decision time $\tau$, and the dispersions $s_v, s_\tau$.

\exonerec{cdm}{CDM}{intercept_only}{intercept-only}
\exoneecdf{cdm}{CDM}{intercept_only}{intercept-only}
\exonerec{cdm}{CDM}{regressed}{regressed}
\exoneecdf{cdm}{CDM}{regressed}{regressed}
\exonerec{cdm}{CDM}{fixed_regressed}{fixed-regressed}
\exoneecdf{cdm}{CDM}{fixed_regressed}{fixed-regressed}
\exonerec{cdm}{CDM}{interaction}{interaction}
\exoneecdf{cdm}{CDM}{interaction}{interaction}
\exonerec{cdm}{CDM}{fixed}{fixed}
\exoneecdf{cdm}{CDM}{fixed}{fixed}
\exonerec{cdm}{CDM}{mixed_predictor}{mixed-predictor}
\exoneecdf{cdm}{CDM}{mixed_predictor}{mixed-predictor}

\FloatBarrier
